\section{Introduction}
\label{sec:intro}

With the advent of large language models (LLMs)~\citep{singh2026openaigpt5card,qiu2026designqwen38nextarchitectureevaluation}, voice assistants have shifted from cascaded pipelines of automatic speech recognition (ASR), text processing, and text-to-speech (TTS)~\citep{lin-etal-2024-advancing} toward end-to-end speech LLMs that directly take speech as input and generate speech as output~\citep{zhang-etal-2023-speechgpt,NEURIPS2024_ecfa5340,ICLR2025_90d1fc07,zeng2024glm4voiceintelligenthumanlikeendtoend}.
Most early systems remain half-duplex: the user and model take turns, and response generation begins only after the user is judged to have finished speaking~\citep{SileroVAD,liao2025flexduopluggableenablingfullduplex,11463929}.
Human conversation, however, includes interruptions, overlap, and backchannels.
This has motivated full-duplex spoken dialog models that continuously listen while generating speech, allowing both sides to speak and respond even when their speech overlaps~\citep{defossez2024moshispeechtextfoundationmodel,11463413,kim2026raonspeechtechnicalreport,cui2026minicpmo45realtimefullduplex}. 

\begin{figure}[t]
    \centering
    \includegraphics[width=0.9\textwidth]{Figures/Figure1.pdf}
    \caption{
    Overview of conventional full-duplex generation and DiffuPlex.
    Conventional models advance one interaction frame at a time, requiring a new backbone computation at every step (left).
    DiffuPlex predicts multiple future frames in a single backbone wake, consumes a retained prefix, and revises only the unplayed future when new user input diverges from the prediction (right).
    }
    \label{fig:1_overview}
    \vskip -15pt
\end{figure}

Despite enabling continuous interaction, full-duplex models remain fundamentally autoregressive.
Some serialize user and assistant speech into a single sequence and update relatively coarse chunks~\citep{veluri-etal-2024-beyond,zhang-etal-2025-omniflatten,cui2026minicpmo45realtimefullduplex,kim2026raonspeechtechnicalreport}, while multi-stream models keep parallel user and assistant channels and update both at every interaction step~\citep{defossez2024moshispeechtextfoundationmodel,11463413,ohashi2026multifacetedinteractivityalignmentfullduplex}.
This finer granularity enables prompt reactions to interruptions, overlaps, and backchannels, but increases backbone computation.
Moshi~\citep{defossez2024moshispeechtextfoundationmodel} and PersonaPlex~\citep{11463413}, for example, process one frame every 80ms, requiring a new temporal-backbone forward pass every 80ms throughout the conversation.

Fine-grained interaction, however, does not necessarily require a fresh backbone computation at every frame.
Across 38h of Switchboard dyadic telephone conversations~\citep{225858}, 95.3\% of active conversational time contains only one speaking participant~\citep{10.3389/fpsyg.2015.00731}, leaving long stretches in which the other side remains silent.
Even during speech, upcoming content is planned beyond individual frames: response preparation can begin before the interlocutor finishes speaking~\citep{10.3389/fpsyg.2017.00393}, and speech planning can extend beyond individual words~\citep{MARTIN2010177}.
These observations suggest that several future interaction frames may be predicted together.
Text generation has exploited similar redundancy through speculative decoding~\citep{pmlr-v202-leviathan23a,chen2023acceleratinglargelanguagemodel} and diffusion language models~\citep{NEURIPS2025_48b383b2,ICLR2025_7ede97c3,fu2026nemotronlabsdiffusiontrimodelanguagemodel}.
However, extending such parallel prediction to full-duplex speech is challenging because the future user speech needed to verify the assistant draft has not yet arrived.

We introduce \textbf{DiffuPlex}, which replaces frame-by-frame autoregressive rollout with masked prediction of a short future block.
At each temporal-backbone forward pass, which we call a \textit{backbone wake}, DiffuPlex predicts a block of future user and assistant frames by replacing unavailable future inputs with learned MASK tokens.
It retains a confident contiguous prefix of the assistant draft and plays it at the original interaction rate while user speech continues to arrive.
The user predictions are stored for consistency checking.
If the observed user input diverges from the prediction, DiffuPlex keeps all assistant frames that have already been played and revises only the remaining future.
Figure~\ref{fig:1_overview} illustrates this process of predicting a future block and revising its unplayed portion as the interaction unfolds.
We refer to this repeated process of predicting, consuming, and revising a masked future as \textit{rolling masked diffusion}.

We evaluate DiffuPlex on full-duplex interaction, spoken-language capability, human evaluation, and inference efficiency.
We consider two inference policies over the same predictor: DiffuPlex-\textsc{Listen} consumes multiple future frames only when they predict assistant silence, whereas DiffuPlex-\textsc{Speak} can also consume predicted assistant speech.
On Full-Duplex-Bench v1.0, the two policies advance the interaction by 1.91 and 2.23 frames per backbone wake on average while remaining close to the autoregressive backbone in interaction behavior.
On the fixed 60-clip runtime subset, they achieve 1.46$\times$ and 1.59$\times$ deployment-path wall-clock speedups and 1.61$\times$ and 1.80$\times$ Core LM speedups, with all measured invocations completing within 80ms.
The more conservative \textsc{Listen} policy also preserves human-rated speech and conversational quality, while extending prediction into assistant speech exposes a speech-quality trade-off that we analyze further.
To test whether speedups can be achieved without future prediction, we compare against training-free baselines that defer autoregressive execution and emit silence on skipped frames, at the cost of degraded full-duplex conversational dynamics.
DiffuPlex instead predicts the intervening interaction, reducing backbone calls while better preserving these dynamics.